% University project report for Graduate Deep RL Paper Chains
\documentclass[12pt]{article}
\usepackage[margin=1in]{geometry}
\usepackage{amsmath,amssymb,amsfonts}
\usepackage{booktabs}
\usepackage{graphicx}
\usepackage{hyperref}
\usepackage{cleveref}
\usepackage{enumitem}
\usepackage{siunitx}
\usepackage{listings}
\usepackage[dvipsnames]{xcolor}
\usepackage{pgfplots}
\pgfplotsset{compat=1.18}
\hypersetup{colorlinks=true,linkcolor=MidnightBlue,citecolor=ForestGreen,urlcolor=blue}

\newcommand{\code}[1]{\texttt{#1}}
\newcommand{\E}{\mathbb{E}}

\title{Graduate Deep Reinforcement Learning Paper Chains\\Comprehensive Project Report}
\author{Your Name \and Partner Name\\Course: Deep Reinforcement Learning\\Institution: <University Name>}
\date{\today}

\lstset{basicstyle=\ttfamily\small,breaklines=true,frame=single,backgroundcolor=\color{gray!5},keywordstyle=\color{blue!70!black},commentstyle=\color{gray!70},stringstyle=\color{red!70!black}}

\begin{document}
\maketitle
\tableofcontents
\listoftables
\listoffigures

\section{Introduction}
This project reproduces and analyzes five canonical deep RL paper chains spanning value-based, policy-gradient/trust-region, actor--critic with maximum entropy, model-based, and multi-agent learning. The goals are: (i) implement faithful yet compact baselines, (ii) evaluate them on standardized benchmarks, (iii) compare theoretical guarantees with empirical behaviour, and (iv) provide reproducible artifacts suitable for graduate coursework.

\section{Related Work}
\begin{itemize}[leftmargin=*]
    \item Tabular Q-learning convergence \cite{watkins1992} and stability refinements (Double DQN \cite{vanhasselt2016double}, Dueling DQN \cite{wang2016dueling}).
    \item Trust-region methods TRPO \cite{schulman2015trpo} and PPO \cite{schulman2017ppo}; constrained policy optimization \cite{achiam2017cpo}.
    \item Actor--critic advances including A3C \cite{mnih2016a3c} and maximum-entropy SAC \cite{haarnoja2018sac}.
    \item Model-based rollouts from Dyna-Q \cite{sutton1990dyna} to ensemble TRPO \cite{kurutach2018metrpo}.
    \item Multi-agent formulations: minimax Q \cite{littman1994}, centralized training with decentralized execution (MADDPG \cite{lowe2017maddpg}), and value factorization (QMIX \cite{rashid2018qmix}).
\end{itemize}

\section{Methods}
\subsection{Codebase Overview}
All code lives in \code{clean\_grad\_rl/}. Entry scripts implement the five chains (see \cref{tab:commands}). Shared utilities handle seeding, environment creation, reproducible logging, and evaluation.

\begin{table}[h]
    \centering
    \caption{Entry-point commands (laptop-friendly defaults).}\label{tab:commands}
    \begin{tabular}{ll}
        \toprule
        Chain & Command \\
        \midrule
        Single run & \code{python scripts/run\_experiment.py --chain value --algo dqn --env CartPole-v1 --steps 20000 --seed 0} \\
        Full suite & \code{python scripts/run\_suite.py --suite configs/suites/fast\_5chain.yaml --mode full} \\
        Smoke suite & \code{python scripts/run\_suite.py --suite configs/suites/fast\_5chain.yaml --mode smoke} \\
        Plot build & \code{python scripts/generate\_plots.py --suite-dir outputs/suite\_reports --runs-root outputs/runs} \\
        Table sync & \code{python scripts/build\_report\_tables.py --aggregate-csv outputs/suite\_reports/tables/aggregate\_metrics.csv --out-root .} \\
        \bottomrule
    \end{tabular}
\end{table}

\subsection{Chain Descriptions}
\paragraph{Value-based.} DQN with replay and target networks; optional Double, Dueling, and QR-DQN variants.
\paragraph{Policy-gradient / trust-region.} REINFORCE baseline; PPO clip surrogate; TRPO with CG and line search; CPO-lite using Lagrangian penalty for safety cost.
\paragraph{Actor--critic / max-entropy.} A2C as synchronous A3C stand-in; SAC with automatic entropy tuning.
\paragraph{Model-based.} Tabular Dyna-Q with planning updates; lightweight ME-TRPO-style ensemble dynamics with periodic TRPO updates using imagined rollouts.
\paragraph{Multi-agent.} QMIX monotonic value factorization and MADDPG centralized critics on PettingZoo MPE simple\_spread.

\section{Experimental Setup}\label{sec:experimental}
\textbf{Hardware:} document CPU/GPU, RAM.\\
\textbf{Software:} Python 3.10+, PyTorch 2.2.2, Gymnasium 0.29.1, PettingZoo MPE (pinned in \code{requirements.txt}).\\
\textbf{Environments:} CartPole/LunarLander, MountainCarContinuous, Pendulum, CliffWalking, MPE simple\_spread.\\
\textbf{Seeding:} \code{--seed} flag; vector environments offset seeds per worker.\\
\textbf{Logging:} JSON metrics in \code{outputs/runs/.../metrics.json}, aggregate CSV in \code{outputs/suite\_reports/tables/aggregate\_metrics.csv}.

\section{Experiments and Results}
Populate the following tables after running the scripts.

\subsection{Value-Based Results}
\IfFileExists{outputs/suite_reports/figures/value_dqn_CartPole-v1.png}{
\begin{figure}[h]
    \centering
    \includegraphics[width=0.82\linewidth]{outputs/suite_reports/figures/value_dqn_CartPole-v1.png}
    \caption{Value-based learning curve generated from suite runs.}
\end{figure}
}{\textit{No generated value figure found yet. Run suite + plot scripts.}}
\begin{table}[h]
\centering
\caption{Value-based results (auto-filled).}
\label{tab:dqnresults}
\begin{tabular}{lccc}
\toprule
Variant & Seeds & Mean Return & Std Across Seeds \\
\midrule
DQN & 3 & 206.07 & 184.89 \\
Rainbow-lite & 3 & -98.07 & 117.07 \\
\bottomrule
\end{tabular}
\end{table}


\subsection{Policy Gradient / Trust Region}
\IfFileExists{outputs/suite_reports/figures/policy_ppo_CartPole-v1.png}{
\begin{figure}[h]
    \centering
    \includegraphics[width=0.82\linewidth]{outputs/suite_reports/figures/policy_ppo_CartPole-v1.png}
    \caption{Policy gradient curve generated from suite runs.}
\end{figure}
}{\textit{No generated policy figure found yet.}}
\begin{table}[h]
\centering
\caption{Policy-gradient / trust-region results.}
\label{tab:policyresults}
\begin{tabular}{lcccc}
\toprule
Algo & Env & Steps & Mean Return & Mean Cost \\
\midrule
REINFORCE & & & & \\
PPO & & & & \\
TRPO & & & & \\
CPO-lite & & & & \\
\bottomrule
\end{tabular}
\end{table}


\subsection{Actor--Critic / Max-Entropy}
\IfFileExists{outputs/suite_reports/figures/actor_critic_sac_Pendulum-v1.png}{
\begin{figure}[h]
    \centering
    \includegraphics[width=0.82\linewidth]{outputs/suite_reports/figures/actor_critic_sac_Pendulum-v1.png}
    \caption{Actor--critic curve generated from suite runs.}
\end{figure}
}{\textit{No generated actor-critic figure found yet.}}
\begin{table}[h]
\centering
\caption{Actor--critic / maximum-entropy results.}
\label{tab:acresults}
\begin{tabular}{lccc}
\toprule
Algo & Steps & Mean Return & Std Return \\
\midrule
A2C (A3C stand-in) & & & \\
SAC & & & \\
\bottomrule
\end{tabular}
\end{table}


\subsection{Model-Based}
\IfFileExists{outputs/suite_reports/figures/model_based_mbpo_lite_Pendulum-v1.png}{
\begin{figure}[h]
    \centering
    \includegraphics[width=0.82\linewidth]{outputs/suite_reports/figures/model_based_mbpo_lite_Pendulum-v1.png}
    \caption{Model-based curve generated from suite runs.}
\end{figure}
}{\textit{No generated model-based figure found yet.}}
\begin{table}[h]
\centering
\caption{Model-based results (auto-filled).}
\label{tab:modelresults}
\begin{tabular}{lccc}
\toprule
Algo & Environment & Seeds & Mean Return \\
\midrule
Dyna-Q & CliffWalking-v0 & 3 & -47.52 \\
MBPO-lite & Pendulum-v1 & 3 & -950.00 \\
\bottomrule
\end{tabular}
\end{table}


\subsection{Multi-Agent}
\IfFileExists{outputs/suite_reports/figures/marl_qmix_lite_simple_spread_v3.png}{
\begin{figure}[h]
    \centering
    \includegraphics[width=0.82\linewidth]{outputs/suite_reports/figures/marl_qmix_lite_simple_spread_v3.png}
    \caption{Multi-agent curve generated from suite runs.}
\end{figure}
}{\textit{No generated MARL figure found yet.}}
\begin{table}[h]
\centering
\caption{Multi-agent results (auto-filled).}
\label{tab:marlresults}
\begin{tabular}{lcc}
\toprule
Algo & Seeds & Episode Reward Mean \\
\midrule
IPPO & 3 & -72.08 \\
QMIX-lite & 3 & -34.58 \\
\bottomrule
\end{tabular}
\end{table}


\section{Ablations}
Suggested ablations: replay buffer size, Double vs. vanilla DQN; PPO clip range; SAC entropy coefficient fixed vs. auto; Dyna-Q planning steps; QMIX worker counts.

\section{Discussion}
Relate empirical findings to theoretical expectations (see \S\ref{sec:experimental} for setup): overestimation bias reduction, KL constraint stability, entropy-driven exploration, model bias vs. rollout horizon, multi-agent non-stationarity.

\section{Limitations and Ethics}
Note approximations (A2C for A3C, simplified CPO), environment realism limits, and responsible reporting of randomness and hyperparameters.

\section{Reproducibility Checklist}
\begin{itemize}[leftmargin=*]
    \item Code and scripts: \code{clean\_grad\_rl/scripts/*.py}
    \item Configs: YAML defaults under \code{clean\_grad\_rl/configs}
    \item Seeds: recorded in logs; set via CLI
    \item Dependencies: \code{requirements.txt}
    \item Hardware/software: state in \S\ref{sec:experimental}
    \item Data: environments auto-downloaded when required
    \item Metrics: JSON/CSV/PNG artifacts in \code{outputs/runs} and \code{outputs/suite\_reports}
\end{itemize}

\section{Conclusion}
Summarize which methods are most stable, sample-efficient, or safe for the tested domains, and outline future extensions (Rainbow, true CPO, PlaNet/MuZero, SMACv2).

\bibliographystyle{plain}
\bibliography{references}

\appendix
\section{Theoretical Notes}
For detailed derivations and proofs, see \code{solutions.tex}; reuse its content here as needed.

\end{document}
